% ====================================================================
% §5 평형 진행률 ξ_eq 와 폭 w_j (N5·N4) — 재설계: logistic 유도 선행 → 폭 모수화
% ====================================================================
\section{평형 진행률 $\xi_\eq$ 와 폭 $w_j$}\label{sec:width}
중심이 정해지면 다음으로 평형 진행률 $\xi_{\eq,j}$ 와 그 폭 $w_j$ 를 정한다. 이 둘이 평형 종(다음 절의 peak)의
모양을 결정한다. 절 순서는 유도가 모수화에 선행하도록 짠다 --- 먼저 평형 등온선을 정$\cdot$역 속도의 균형에서
logistic 으로 닫아 자연 폭 척도 $RT/F$ 를 드러내고(\S\ref{sec:width-logistic}), 그 위에서 폭 다중도 $n_j$ 로
일반화한 뒤 폭의 이중지위를 세운다(\S\ref{sec:width-w}). 이 운동학 경로는 Part 0(\S\ref{sec:sm-electro})의 평형
통계 경로와 독립이며, 두 경로가 같은 logistic 에 닿는 것이 \S\ref{sec:dist} 의 통합이다.

\subsection{평형 진행률 $\xi_\eq$ --- logistic 의 유도}\label{sec:width-logistic}
평형 등온선은 \S\ref{sec:center} 의 평형 조건을 정$\cdot$역 속도의 균형으로 회수하면 logistic 으로 닫힌다.
\textbf{(a) 출발 --- Eyring 속도식과 비대칭 장벽.} 활성화 장벽 $\Delta G_a$ 이상을 갖출 확률은 Boltzmann 분포
$\propto e^{-\Delta G_a/RT}$ 이고, 전이상태 이론의 시도 빈도 $k_0=k_BT/h$ 를 곱하면 Eyring 속도식
$k\simeq k_0e^{-\Delta G_a/RT}$ 다(그림~\ref{fig:barrier})\cite{eyring1935}. 구동력 $\mathcal A_j=sF(V_n-U_j)$ 가 정$\cdot$역 장벽을
비대칭 이동시켜(분율 위치 $\chi$ --- 비평형 열역학 기반 전하전달 동역학의 표준 분해\cite{bazant2013}), 자리당 정$\cdot$역 속도가
\begin{equation}
r_j^+=k_0\,e^{-(\Delta G_{a,j}-\chi\mathcal A_j)/RT},\qquad
r_j^-=k_0\,e^{-(\Delta G_{a,j}+(1-\chi)\mathcal A_j)/RT}.
\label{eq:bv}
\end{equation}
\textbf{(b) 연산 --- 비를 취해 detailed balance.} 두 속도의 비를 취하면 $k_0,e^{-\Delta G_a/RT}$ 가 약분되고 지수의
$\chi$ 가 상쇄되어($\chi\mathcal A+(1-\chi)\mathcal A=\mathcal A$)
\begin{equation}
\frac{r_j^+}{r_j^-}=\exp\!\Big[\frac{\chi\mathcal A_j+(1-\chi)\mathcal A_j}{RT}\Big]=\exp\!\Big[\frac{\mathcal A_j}{RT}\Big]
\label{eq:db}
\end{equation}
--- 평형비가 $\chi$ 와 무관한 Boltzmann 비라는 detailed balance 다($\chi+(1-\chi)=1$ 합-1 제약이 강제한 결과). 운동
방정식은 정방향이 찬 자리($1-\xi$)에서, 역방향이 빈 자리($\xi$)에서 출발하므로
$\dd\xi/\dd t=r^+(1-\xi)-r^-\xi=k_j(\xi_\eq-\xi)$($k_j\equiv r^++r^-$). \textbf{(c) 중간식 --- 정지점에서 logit.}
정지점 $\dd\xi/\dd t=0$ 에서 $r^+(1-\xi_\eq)=r^-\xi_\eq$(정$\cdot$역 플럭스의 교점 --- 그림~\ref{fig:flux}), 곧
비 식~\eqref{eq:db} 로
\begin{equation}
\frac{\xi_\eq}{1-\xi_\eq}=e^{\mathcal A_j/RT}=e^{sF(V_n-U_j)/RT}
\quad\Longrightarrow\quad
\xi_\eq=\frac{1}{1+e^{-sF(V_n-U_j)/RT}}
\label{eq:logisticsolve}
\end{equation}
(둘째 식은 분모$\cdot$분자를 $e^{\mathcal A/RT}$ 로 나눈 것). 이 bare logistic 의 중심 기울기는
$\dd\xi_\eq/\dd V|_{1/2}=sF/(4RT)$ 로, 폭 척도 $RT/F$ 를 정의한다 --- 열역학은 detailed balance~\eqref{eq:db}
한 곳으로만 들어왔고, logistic 은 속도식의 정지점에서 \emph{유도된} 평형이다. 다음 소절이 이 폭 척도를 다중도
$n_j$ 로 일반화해 결과 박스를 닫는다. 본문이 여기서 빌려 쓴 두 양 --- 시도 빈도 $k_0=k_BT/h$ 와
\S\ref{sec:lag} 가 $\Delta H_a-T\Delta S_a$ 로 가를 활성화 자유에너지 --- 의 미시 기원은 아래 배경 박스가
전이상태 이론의 분배함수로 닫는다.

\begin{bgbox}
\textbf{전이상태 이론의 분배함수와 활성화 엔트로피 ---} 위 (a) 의 시도 빈도 $k_0=k_BT/h$ 와,
\S\ref{sec:lag} 이 쓰는 활성화 엔트로피 $\Delta S_a$ 는 전이상태 이론(transition-state theory, TST)의
분배함수에서 유도된다 --- 이 박스는 그 미시 기원만 자족적으로 정리한다(본문 흐름은 이 박스 없이도 읽히고,
\S\ref{sec:lag} 은 여기서 세운 갈래를 \emph{쓴다}).

\textbf{(a) 출발 --- 두 전제.} TST 는 두 가정 위에 선다: (i) \emph{안장점 준평형} --- 안장점의
활성복합체(activated complex)가 반응물과 Boltzmann 평형을 이룬다; (ii) \emph{재교차 무시} --- 안장점을
정방향으로 넘는 궤적은 되돌지 않고 모두 생성물이 된다(투과계수 $\kappa=1$)\cite{laidlerking1983}. 반응좌표
모드는 길이 $\delta$ 구간의 자유 1D 병진으로 다룬다 --- 변분 TST·터널링 보정은 범위 밖이다.

\textbf{(b) 연산 --- $k_BT/h$ 의 미시 기원.} 안장점에서 반응좌표 모드를 뗀 활성복합체 분배함수를
$q^{\ddagger}$, 반응물의 것을 $q_R$ 로 둔다. 뗀 반응좌표 모드는 유효질량 $m^{*}$ 의 1D 병진 분배함수
\begin{equation}
q_\mathrm{rc}=\frac{(2\pi m^{*}k_BT)^{1/2}\,\delta}{h}
\label{eq:tst-qrc}
\end{equation}
--- Part 0 의 $q(T)$~\eqref{eq:partfn} 와 같은 위상공간/$h$ 눈금으로 센 상태 수 --- 이고, 그 복합체가
안장점을 정방향으로 넘는 평균 속도는 반쪽 Maxwell 평균 $\langle v\rangle=(k_BT/2\pi m^{*})^{1/2}$ 다. 길이
$\delta$ 에 갇힌 복합체의 통과 빈도 $\langle v\rangle/\delta$ 에 식~\eqref{eq:tst-qrc} 를 곱하면 $\delta$ 와
$m^{*}$ 가 정확히 상쇄되어
\begin{equation}
\frac{\langle v\rangle}{\delta}\,q_\mathrm{rc}
=\frac{1}{\delta}\Big(\frac{k_BT}{2\pi m^{*}}\Big)^{1/2}\frac{(2\pi m^{*}k_BT)^{1/2}\delta}{h}
=\frac{k_BT}{h}
\label{eq:tst-freq}
\end{equation}
만 남는다 --- $k_0=k_BT/h$ 는 임의로 고른 상수가 아니라 유효질량·구간 길이에 무관한 \emph{장벽 통과의 보편
빈도}다($298$ K 에서 $6.2\times10^{12}$ /s).

\textbf{(c) 중간식 --- 준평형 인구와 자유에너지 묶기.} 전제 (i) 로 활성복합체 대 반응물의 인구비가 분배함수
비에 장벽 바닥차 $\Delta E_0$ 의 Boltzmann 인자를 곱한 값이므로, 통과 빈도에 이 인구를 실은 속도는
\begin{equation}
k=\frac{k_BT}{h}\,\frac{q^{\ddagger}}{q_R}\,e^{-\Delta E_0/RT}
\label{eq:tst-rate}
\end{equation}
다\cite{glasstone1941,mcquarrie1976}. 본문 식~\eqref{eq:bv} 의 활성화 자유에너지 $\Delta G_a$ 를 분배함수로
쓰면 --- 재정의가 아니라 같은 양의 미시 표기 \emph{식별}이다 ---
\begin{equation}
\Delta G_a=-RT\ln\frac{q^{\ddagger}}{q_R}+\Delta E_0
\label{eq:tst-dG}
\end{equation}
이고, 식~\eqref{eq:tst-rate} 는 $k=(k_BT/h)\,e^{-\Delta G_a/RT}$ 로 정리되어 식~\eqref{eq:bv} 가 쓰는 Eyring
형과 같아진다.

\textbf{(d) 박스 --- 활성화 엔트로피의 미시 의미.} $\Delta G_a=\Delta H_a-T\Delta S_a$ 로 가르면
\begin{equation}
\boxed{\;k=\frac{k_BT}{h}\,e^{\Delta S_a/R}\,e^{-\Delta H_a/RT},
\qquad \Delta S_a=R\ln\frac{q^{\ddagger}}{q_R}\;.}
\label{eq:tst-box}
\end{equation}
활성화 엔트로피는 (반응좌표를 뗀) 활성복합체 분배함수 $q^{\ddagger}$ 와 반응물 분배함수 $q_R$ 의 비의 로그다 --- 활성복합체의
자유도가 반응물보다 \emph{조이면}($q^{\ddagger}<q_R$) $\Delta S_a<0$ 라 prefactor 가 눌리고, \emph{풀리면}
($q^{\ddagger}>q_R$) $\Delta S_a>0$ 라 커진다.

\textbf{Part 0 접속.} 이 ``로그-분배함수 $=$ 엔트로피'' 대응은 Part 0 의 내부 자유도 언어와 같다 --- Part 0 은
단일 자리 내부 분배함수 $q(T)$(식~\eqref{eq:partfn})에 연산 $s_\mathrm{int}=k_B\,\partial(T\ln q)/\partial T$
(식~\eqref{eq:sm-sint})를 걸어 자리 엔트로피를 읽었고, 여기서는 같은 연산을 활성화의 비 $q^{\ddagger}/q_R$ 에
걸어 활성화 엔트로피를 읽는다(비의 잔여 $T$-의존은 $\Delta H_a$ 쪽으로 넘기는 표준 읽기). \emph{가드} ---
여기 $\Delta S_a$ 는 \emph{활성화}(비가역 동역학, 안장점 대 반응물의 비)이고, 전이 중심의
$\Delta S_\rxn$(\S\ref{sec:center} --- 반응·평형, 생성물 대 반응물의 비)과 차원은 같으나 물리가
다르다.\footnote{여기 $q^{\ddagger},q_R$(TST 분배함수)와 Part 0 의 $q(T)$(자리 내부 분배함수)는 모두 통계역학
분배함수로, N6--N9 의 용량 좌표 $q{=}Q/Q_\cell$ 와는 무관하다(식~\eqref{eq:partfn} 각주와 같은 가드).}
\begin{verifybox}
검산 --- (i) 차원: $k_BT/h$ 는 에너지/(에너지·시간)$=$[1/s]이고 두 지수는 무차원이라 $k$ 는 빈도
차원이다(식~\eqref{eq:tst-freq}) --- \S\ref{sec:lag} 완화율 $k_j$ 와 같은 눈금. (ii) 고온 극한: 진동 모드의
$q\to k_BT/\hbar\omega$ 라 Eyring 전치인자가
$(k_BT/h)\,(q^{\ddagger}/q_R)\to\prod_{i=1}^{N}\nu_i^{R}\big/\prod_{i=1}^{N-1}\nu_i^{\ddagger}$
--- 온도$\cdot$유효질량에 무관한 \emph{고전 시도빈도}(진동수$\cdot$곡률 비) --- 로 유한 수렴한다
($q_R$ 는 반응물 전 모드$\cdot$$q^{\ddagger}$ 는 반응좌표를 뺀 모드라 비 자체는 $\propto1/T$;
유한 수렴의 주체는 비가 아니라 $k_BT/h$ 와 묶인 전치인자다). (iii) 원형 회수: 활성화가 (반응좌표를 뺀 비교에서) 자유도를 조이지도 풀지도 않아 $q^{\ddagger}/q_R=1$ 이면
$\Delta S_a=0$ 이라 식~\eqref{eq:tst-box} 가 $k=(k_BT/h)e^{-\Delta H_a/RT}$ 의 순수 Arrhenius 형으로 환원된다
--- \S\ref{sec:sum} 의 기본값 $\Delta S_a{=}0$ 이 이 극한이다. (iv) 부호:
$q^{\ddagger}\lessgtr q_R\Rightarrow\Delta S_a\lessgtr0\Rightarrow e^{\Delta S_a/R}\lessgtr1$.
\end{verifybox}
\end{bgbox}

\begin{srcbox}
\textbf{다리 --- eq:bv 의 $\chi$ 분할이 어느 정식화인가.} 식~\eqref{eq:bv} 가 구동력
$\mathcal A_j=sF(V_n-U_j)$ 를 정$\cdot$역 장벽에 $\chi:(1-\chi)$ 로 실어 비대칭 이동시키는 것은
Bazant 의 비평형 열역학 전하전달 이론\cite{bazant2013} 이 세운 표준 분해다. 대응은 다음과 같다(왼쪽이
본문 기호, 오른쪽이 원 논문 량):
\begin{center}\footnotesize
\begin{tabular}{@{}ll@{}}
\hline
본문(식~\eqref{eq:bv},~\eqref{eq:db}) & Bazant\cite{bazant2013} 대응량 \\
\hline
분율 위치 $\chi$ / $(1-\chi)$ & 전하전달 대칭인자 $\alpha_c{=}\alpha$ / $\alpha_a{=}1{-}\alpha$ (식[27]) \\
구동력 $\mathcal A_j=sF(V_n-U_j)$ & 전기화학 과전위 $ne\,\eta=\Delta\mu$, $\eta{=}\Delta\phi{-}\Delta\phi^{\eq}$ (식[26]) \\
공통 인자 $k_0e^{-\Delta G_a/RT}$ & 교환 전류 $I_0/ne$ --- 활동도$\cdot$전이상태 구조 포함 (식[29]) \\
평형비 $r_j^+/r_j^-=e^{\mathcal A_j/RT}$ (식~\eqref{eq:db}) & $\eta{=}0$ 에서 알짜 전류 $0$ (식[27]) \\
\hline
\end{tabular}
\end{center}
등치의 핵은 지수 인자다 --- $n{=}1$ 에서 $F{=}N_A e$, $R{=}N_A\kB$ 이므로
\begin{equation}
\frac{\mathcal A_j}{RT}=\frac{sF(V_n-U_j)}{RT}=\frac{e\,[\,s(V_n-U_j)\,]}{\kB T}
\;\;\longleftrightarrow\;\;\frac{ne\,\eta}{\kB T},\qquad \eta \leftrightarrow s(V_n-U_j),
\label{eq:br-bazant2013-1}
\end{equation}
곧 우리 부호 구동전압 $s(V_n-U_j)$ 가 원 논문 과전위 $\eta$ 에 대응한다(이 $\eta$ 는 원 논문 기호 --- \S\ref{sec:broadening} 의 겉보기-중심 국소 과전위 $\eta$ 와 별개다). \emph{방법 요지} --- Bazant 는 알짜
속도를 전이상태 excess 화학퍼텐셜 기준 반응물$\cdot$생성물 화학퍼텐셜의 지수차(식[7])로 적고, 구동력을
$\eta=\Delta\mu/ne$(식[26])로 특수화한 뒤 장벽을 $\alpha$ 로 비대칭 분할해 일반화 Butler--Volmer(식[27])를
얻는다. \emph{가정 차} --- 원 논문의 교환전류 $I_0$(식[29])는 농축용액 활동도계수와 전이상태 활동도계수
$\gamma_\ddagger$ 를 실어 전치인자가 \emph{조성 의존}이나, 식~\eqref{eq:bv} 는 이를 조성무관 상수
$k_0e^{-\Delta G_a/RT}$ 로 흡수하고(이상 교환속도) 조성 의존을 전부 구동력 $\mathcal A_j$ 와 logistic 로만
받는다.
\end{srcbox}
% 자산(v1.0.21 신규): [V21-Q3-01 반응좌표 1D 분배함수 eq:tst-qrc] [V21-Q3-02 보편 빈도 k_BT/h 유도 eq:tst-freq]
%   [V21-Q3-03 TST 속도식 eq:tst-rate] [V21-Q3-04 ΔG_a 미시 식별 eq:tst-dG]
%   [V21-Q3-05 ΔS_a=R ln(q‡/q_R) boxed eq:tst-box + ΔS_rxn 구분 가드·용량 q 가드]

\begin{figure}[t]
\centering
% T4 fig:barrier  (quantitative-density)
% Eyring activation landscape (eq:bv), reaction coordinate x, wells at x=1,5, TS at x=3.
% G_base(x)=0.1+0.9 sin^2[pi(x-1)/4] on [1,5]; driven G(x)=G_base-A*phi(x), phi=(x-1)/4.
% A=0.5, chi=1/2 -> forward barrier dG_a-chi A=0.65, reverse dG_a+(1-chi)A=1.15, peak down chi A=0.25.
\begin{tikzpicture}[x=0.70cm,y=2.4cm,>=Latex]
% =============== (a) equilibrium ===============
\begin{scope}
\draw[->] (-0.3,-0.5) -- (6.5,-0.5) node[right,font=\scriptsize] {rxn coord};
\draw[->] (-0.3,-0.5) -- (-0.3,1.18) node[above,font=\scriptsize] {free energy $G$};
\draw[thick] plot[smooth] coordinates {(0,0.1) (0.5,0.1) (1,0.1) (1.5,0.2318) (2,0.55) (2.5,0.8682) (3,1.0) (3.5,0.8682) (4,0.55) (4.5,0.2318) (5,0.1) (5.5,0.1) (6,0.1)};
\node[below,font=\scriptsize] at (1,0.06) {filled site};
\node[below,font=\scriptsize] at (5,0.06) {empty site $+$ ion};
\node[above,font=\scriptsize] at (3,1.0) {transition state};
% dG_a arrow at x=2
\draw[densely dotted,black!55] (1,0.1) -- (2.05,0.1);
\draw[densely dotted,black!55] (3,1.0) -- (2.05,1.0);
\draw[<->] (2.05,0.1) -- (2.05,1.0) node[pos=0.5,left,font=\scriptsize] {$\Delta G_a{=}0.90$};
\node[font=\scriptsize,below] at (3,-0.68) {(a) equilibrium: $\mathcal A{=}0$ (two wells equal)};
\end{scope}
% =============== (b) driven ===============
\begin{scope}[xshift=6.3cm]
\draw[->] (-0.3,-0.5) -- (6.7,-0.5) node[right,font=\scriptsize] {rxn coord};
% reference (a) dotted
\draw[densely dotted,thin,black!50] plot[smooth] coordinates {(0,0.1) (1,0.1) (1.5,0.2318) (2,0.55) (2.5,0.8682) (3,1.0) (3.5,0.8682) (4,0.55) (4.5,0.2318) (5,0.1) (6,0.1)};
% driven curve
\draw[thick] plot[smooth] coordinates {(0,0.1) (0.5,0.1) (1,0.1) (1.5,0.1693) (2,0.425) (2.5,0.6807) (3,0.75) (3.5,0.5557) (4,0.175) (4.5,-0.2057) (5,-0.4) (5.5,-0.4) (6,-0.4)};
% peak down by chi A
\draw[->,thick] (3,0.99) -- (3,0.77);
\node[above,font=\scriptsize] at (3,1.0) {peak $\downarrow\,\chi\mathcal A{=}0.25$};
% forward barrier at x=1.7
\draw[densely dotted,black!55] (1,0.1) -- (1.7,0.1);
\draw[densely dotted,black!55] (3,0.75) -- (1.7,0.75);
\draw[<->] (1.7,0.1) -- (1.7,0.75) node[pos=0.5,left,font=\scriptsize,align=right] {$\Delta G_a{-}\chi\mathcal A$\\$=0.65$};
% reverse barrier at x=4.3
\draw[densely dotted,black!55] (5,-0.4) -- (4.3,-0.4);
\draw[densely dotted,black!55] (3,0.75) -- (4.3,0.75);
\draw[<->] (4.3,-0.4) -- (4.3,0.75) node[pos=0.5,right,font=\scriptsize,align=left] {$\Delta G_a{+}(1{-}\chi)\mathcal A$\\$=1.15$};
% affinity A drop at x=5.8
\draw[densely dotted,black!55] (5,0.1) -- (5.85,0.1);
\draw[densely dotted,black!55] (5,-0.4) -- (5.85,-0.4);
\draw[<->] (5.85,0.1) -- (5.85,-0.4) node[pos=0.5,right,font=\scriptsize] {$\mathcal A{=}0.50$};
\node[font=\scriptsize,below] at (3,-0.68) {(b) driven $\mathcal A{>}0$ ($\chi{=}\tfrac12$)};
\node[font=\scriptsize,anchor=south west] at (-0.2,-0.48) {$\Delta H_a^\eff{=}\Delta H_a{-}\chi_d\Omega$};
\end{scope}
\end{tikzpicture}
\caption{활성화 장벽 그림 --- Eyring 속도식(식~\eqref{eq:bv})의 반응 좌표 개형(좌표는 모형 이중웰 퍼텐셜
$G_\mathrm{base}{=}0.1{+}0.9\sin^2[\pi(x{-}1)/4]$ 의 수치 평가, 구동은 진행률 $\phi$ 를 따라 $\mathcal A$ 만큼 기울임).
(a) 평형($\mathcal A{=}0$) --- 두 골 같은 높이, 장벽 $\Delta G_a{=}0.90$. (b) 구동력 $\mathcal A_j{=}sF(V_n{-}U_j){>}0$
($\chi{=}\tfrac12$)가 도착 골을 $\mathcal A{=}0.50$ 내려 정점을 $\chi\mathcal A{=}0.25$ 낮추면, 정방향 장벽
$\Delta G_a{-}\chi\mathcal A{=}0.65$ 로 하강$\cdot$역방향 $\Delta G_a{+}(1{-}\chi)\mathcal A{=}1.15$ 로 상승(점선 $=$(a)).
$\chi\!\cdot\!(1{-}\chi)$ 분할이 detailed balance(식~\eqref{eq:db})를 강제하고, 유효 장벽은 $\Delta H_a^\eff{=}\Delta H_a{-}\chi_d\Omega$
다. 이 그림이 logistic(식~\eqref{eq:xieq})과 $L_q$ 의 공통 출발점이다.}
\label{fig:barrier}
\end{figure}

\begin{figure}[t]
\centering
% T5 fig:flux  (quantitative-density)
% Forward flux r+(1-xi) vs reverse r- xi; stationary point = intersection = xi_eq (eq:logisticsolve).
% r- = 1 fixed => reverse line y=xi common to all cases; three affinities overlaid:
%   A/RT=0     r+/r-=1  xi_eq=1/2   ; A/RT=ln2  r+/r-=2  xi_eq=2/3 ; A/RT=2 ln2  r+/r-=4  xi_eq=4/5.
\begin{tikzpicture}[x=6.0cm,y=1.42cm,>=Latex]
% ---- readout above the plot (clear of all curves) ----
\node[anchor=south west,align=left,font=\scriptsize] at (-0.02,4.62)
  {reverse $r^-\xi$ ($r^-{=}1$)\ $\cdot$\ forward $r^+(1-\xi)$, $r^+{=}e^{\mathcal A/RT}{=}1,2,4$\\[1pt]
   $\xi_\eq=\dfrac{r^+}{r^++r^-}=\dfrac{1}{1+e^{-\mathcal A/RT}}=\tfrac12,\ \tfrac23,\ \tfrac45$ (교점, on line $y{=}\xi$)};
% ---- axes ----
\draw[->] (0,0) -- (0,4.5) node[left,font=\scriptsize] {flux [arb.]};
\draw[->] (-0.01,0) -- (1.14,0) node[below,font=\scriptsize] {$\xi$};
% ---- common reverse line  y = r- xi ----
\draw[line width=1.3pt] (0,0) -- (1,1);
\node[anchor=south east,font=\scriptsize] at (0.99,0.99) {$r^-\xi$};
% ---- forward lines r+(1-xi) ----
\draw[thick]               (0,4) -- (1,0);
\draw[dashed]              (0,2) -- (1,0);
\draw[gray,densely dashed] (0,1) -- (1,0);
% intercept labels + y ticks
\node[anchor=west,font=\scriptsize]      at (0.02,4.13) {$r^+{=}4$};
\node[anchor=west,font=\scriptsize]      at (0.02,2.13) {$r^+{=}2$};
\node[gray,anchor=west,font=\scriptsize] at (0.02,1.13) {$r^+{=}1$};
\foreach \yv in {1,2,4} {\draw (0.008,\yv) -- (-0.008,\yv) node[left,font=\scriptsize] {\yv};}
% ---- intersections + droplines + x ticks (xi_eq) ----
\foreach \xe in {0.5,0.6667,0.8} {
  \draw[densely dotted,black!55] (\xe,\xe) -- (\xe,0);
  \fill (\xe,\xe) circle (1.1pt);
}
\draw (0.5,0.03) -- (0.5,-0.03) node[below,font=\scriptsize] {$\tfrac12$};
\draw (0.6667,0.03) -- (0.6667,-0.03);
\node[below,font=\scriptsize] at (0.6667,-0.19) {$\tfrac23$};
\draw (0.8,0.03) -- (0.8,-0.03) node[below,font=\scriptsize] {$\tfrac45$};
\draw (1,0.03) -- (1,-0.03) node[below,font=\scriptsize] {1};
% ---- forward family label (clear top-left triangle) ----
\node[anchor=west,font=\scriptsize] at (0.30,3.15) {forward $r^+(1-\xi)$};
% ---- moves-right note (below ticks) ----
\node[anchor=north,font=\scriptsize] at (0.66,-0.42) {$\xi_\eq$ moves right as $\mathcal A\!\uparrow$};
\end{tikzpicture}
\caption{정지점 유도 --- 정$\cdot$역 플럭스의 교점(식~\eqref{eq:logisticsolve}; 좌표는 식 그대로의 수치 평가).
정방향 $r^+(1-\xi)$(감소 직선)와 역방향 $r^-\xi$(공통 증가 직선, $r^-{=}1$)의 교점이 정지점 $\xi_\eq$. affinity 를
$\mathcal A/RT=0,\ln2,2\ln2$ 로 올리면 $r^+/r^-=e^{\mathcal A/RT}=1,2,4$ 라 교점이 $\xi_\eq=\tfrac12\to\tfrac23\to\tfrac45$
로 이동한다(모두 $y{=}\xi$ 선 위) --- 곧 $\xi_\eq=r^+/(r^++r^-)=1/(1+e^{-\mathcal A/RT})$ 로, detailed balance
(식~\eqref{eq:db})가 $\chi$ 와 무관하게 평형을 정한다.}
\label{fig:flux}
\end{figure}

\subsection{폭 $w_j$ --- 이상 극한$\cdot$일반화$\cdot$그 이중지위}\label{sec:width-w}
\textbf{폭 척도의 정의와 다중도 일반화.} 앞 소절 bare logistic(식~\eqref{eq:logisticsolve})의 중심 기울기
$\dd\xi_\eq/\dd V|_{1/2}=sF/(4RT)$ 가 폭 척도 $RT/F$ 를 정의한다. 폭 다중도 $n_j$ 로 일반화하면
\begin{equation}
w_j=\frac{n_jRT}{F}\qquad(\text{기본 폭}),
\label{eq:wbase}
\end{equation}
폭이 직접 주어지면(대신 $w$ 값 자체가 입력되면) $n_j=w_jF/(RT)$ 로 역산해 같은 식에 넣는다.
다중도를 온도 상수가 아니라 온도 함수 $n_j(T)$ 로 두면 폭 $w_j=n_j(T)RT/F$ 가 열적 스케일 위에 잔여 $T$-의존을
얹으며, 그 $\partial w_j/\partial T=(R/F)(n_j(T)+T\,\dd n_j/\dd T)$ 가 가역 발열 config 항에 자기정합 전파된다(Part T
겹침 가중 파생 A(\S\ref{ssec:overlap})에서 유도, 발열 종합식(\S\ref{sec:synthesis})이 받음; 상수 $n_j$ 는 $\dd n_j/\dd T{=}0$ 특수형). 폭을 $n_j$(선택적으로 $n_j(T)$)로 피팅할지는 데이터가 정한다.
폭은 항상 이 한 식을 쓴다(상호작용에 따라 폭을 인위로 좁히는 별도 항은 두지 않는다 --- 그 까닭은 바로 아래
이중지위와 \S\ref{sec:broadening} 가 설명한다).

\textbf{박스 --- 폭 다중도$\cdot$방향 부호 일반화.} 폭 척도 $w_j=n_jRT/F$ 로 묶고 방향 부호 $\sigma_d$ 와 분기
중심 $U_j^{\,d}$(\S\ref{sec:hys})로 식~\eqref{eq:logisticsolve} 를 일반화하면 --- 충전의 진행률은 같은 분포의
여집합 좌표 $\theta=1-\xi$ 라 $s\to\sigma_d$ 치환은 여집합 교환과 동치다(\S\ref{sec:dist}; 전극-중립
일반화는 Chapter 2 \S\ref{sec:lco-direction}) ---
\begin{equation}
\boxed{\;\xi_{\eq,j}(V,T)=\frac{1}{1+\exp\!\big[-\sigma_d(V-U_j^{\,d})/w_j\big]}\;.}
\label{eq:xieq}
\end{equation}
방향 부호가 logistic 인자의 부호로 들어가는 것이 식~\eqref{eq:xieq} 의 $z=\sigma_d(V-U_j^{\,d})/w_j$ 이며,
오버플로를 피하기 위해 $z\ge0/z<0$ 두 등가 형태로 나눠 평가한다(수학적으로
동일)\footnote{여기 무차원 logistic 인자 $z$ 는 이 절 국소 표기이며, \S\ref{sec:sm-mf} 의 배위수(coordination
number) $z$ 나 화학종의 전하수(charge number) $z$ 와는 무관하다.}. 여기서 평가 전위는 \S\ref{sec:pol} 의 내부 전위 $V_n$ 이다. 중심 $U_j^{\,d}$
는 식~\eqref{eq:center} 의 분기 중심(분기 있으면 분기 중심, 없으면 $U_j$)이다.

\begin{signbox}
\textbf{부호.} 방전 $\sigma_d=+1$: $V\!\uparrow\Rightarrow z\!\uparrow\Rightarrow\xi_\eq\!\uparrow$ --- 전위를 올리면
탈리튬화 진행이 는다(규약 일치). 충전 $\sigma_d=-1$: $V\!\uparrow\Rightarrow z\!\downarrow\Rightarrow\xi_\eq\!\downarrow$.
극한 --- $V\gg U_j^{\,d}\Rightarrow\xi_\eq\to1$(방전), $V\ll U_j^{\,d}\Rightarrow0$, $V=U_j^{\,d}\Rightarrow\tfrac12$
(전이 중심). $w_j=RT/F=25.7$ mV($25^\circ$C).
\end{signbox}

\textbf{★폭 $w_j$ 의 이중지위 --- 같은 식, 다른 지위.} 식~\eqref{eq:wbase} 의 \emph{값}은 한 줄이지만, 그것이
\emph{무엇을 뜻하는지}는 전이의 상분리 여부로 갈린다. \emph{단상($\Omega_j\le2RT$, 연속 고용체)}에서는 $w_j=n_jRT/F$
가 등온선 폭에 대한 \emph{검증 가능한 평형 예측}이다 --- 평형 단일 입자 dQ/dV 가 logistic
미분 종(\S\ref{sec:eqpeak} 에서 유도)이고, 그 폭이 곧 $n_jRT/F$ 라($n_j\!\approx\!1$ 이면 $RT/F\!\approx\!25.7$ mV), 측정 폭으로 직접 검증된다. 엄밀히 이 닫힌 폭은 이상 극한($\Omega_j=0$) 기준이고,
$0<\Omega_j<2RT$ 에서는 평형 등온선의 중심 기울기가 $(1-\Omega_j/2RT)$ 배 좁아진 유효 폭을 주되 그
값 역시 평형이 결정하므로 첫째 지위(평형 예측)라는 성격은 변하지 않는다(단상 등온선의 \emph{해석적 성질}
서술이다 --- 모델 폭 항은 여전히 식~\eqref{eq:wbase} 하나이고, 두-상 폭을 이 인자로 좁히는 유효-폭
축소식을 쓰지 않는다는 Part T 파생 C(\S\ref{ssec:weff}) warnbox 의 금지와 무모순$\cdot$동일 취지).
\emph{두-상($\Omega_j>2RT$, 상분리 plateau)}에서는 평형이 예측하는 것이 폭이 아니라 \emph{날카로운 선}(델타)이므로,
같은 $w_j$ 가 더는 평형 예측이 아니라 \emph{broadening 이 정하는 현상학적 자유 피팅 폭}이다(\S\ref{sec:broadening}).
어느 지위인지는 인위적 구분이 아니라 그 전이의 \emph{$\Omega_j$ 값}이 가르며, 같은 식~\eqref{eq:wbase} 가 두 경우 모두에
쓰인다.

이 지위 구분의 흑연 staging 적용은 이렇다. 표~\ref{tab:staging} 의 네 전이는 \emph{초기값} $\Omega_j$ 가 모두 $2RT(\!\approx\!4958$ J/mol
$@298.15$ K$)$ 보다 커 출발점에선 전부 둘째(두-상) 쪽에 놓이지만, 이 $\Omega_j$ 는 거친 초기 추정이라 피팅으로
override 되며, 피팅 후 어느 전이가 어느 지위인지는 그 전이의 \emph{$\Omega_j$ 값}이 가른다 --- 어느 전이가
연속 고용체이고 어느 둘만 두-상 plateau 인가라는 \emph{전이별 분류의 물리 근거}는 \S\ref{sec:broadening-class} 가
세우므로, 이 소절에서 전이 실명을 되풀이하지 않는다. 폭 쪽으로는, 네 전이의 폭 출발값(0.020/0.016/0.014/0.012 V)도 위 $\Omega_j$ 와 같은 자격 --- 피팅이
override 하는 거친 초기값 --- 이고($w$ 폴백 --- $n_j$ 입력을
배제해야 활성화되며 현재 출발은 $n_j{=}1$: \S\ref{sec:broadening} ②), 단상의 $w$ 는 평형이
정하고 두-상의 $w$ 는 피팅이 정한다. 두-상 델타가 실측 종으로 펴지는 세 출처도 \S\ref{sec:broadening} 이 닫는다 --- 이 소절은 폭의 \emph{지위}만 세운다.

그림~\ref{fig:logistic} 가 $\xi_\eq$ 와 그 미분(평형 종)을 함께 보이며, \S\ref{sec:eqpeak}(N6)이 이 종이 평형 peak 임을 닫는다.

\begin{figure}[t]
\centering
% T6 fig:logistic  (quantitative-density)
% Dual axis: left = xi_eq logistic (eq:xieq), right = derivative bell w|dxi/dV|->dxi/dV in 1/V
% (eq:belliden).  T-dependence of width w=n RT/F, n=1, at T=268/298/328 K:
%   w = 23.094 / 25.680 / 28.265 mV ; peak dxi/dV = 1/(4w) = 10.825 / 9.735 / 8.845 /V ;
%   FWHM = 3.5255 w = 81.4 / 90.5 / 99.7 mV.  Center tangent slope of xi_eq = 1/(4w).
\begin{tikzpicture}[x=0.042cm,y=1cm,>=Latex]
% ---- frame + axes ----
\draw[->] (-85,0) -- (92,0);
\node[font=\scriptsize] at (8,-0.60) {$(V-U_j^{\,d})$ [mV]};
\draw[->] (-85,0) -- (-85,3.25) node[above,font=\scriptsize] {$\xi_\eq$};
\draw[->] (85,0) -- (85,3.25) node[above,font=\scriptsize] {$\dd\xi/\dd V$ [1/V]};
% left ticks (xi): 0,0.5,1 -> 0,1.5,3.0
\foreach \yv/\lab in {0/0, 1.5/0.5, 3.0/1} {\draw (-85,\yv) -- (-89,\yv) node[left,font=\scriptsize] {\lab};}
% right ticks (dxi/dV): 0,3,6,9 -> 0,0.818,1.636,2.455
\foreach \yv/\lab in {0.818/3, 1.636/6, 2.455/9} {\draw (85,\yv) -- (89,\yv) node[right,font=\scriptsize] {\lab};}
% x ticks
\foreach \xx in {-80,-40,0,40,80} {\draw (\xx,0.04) -- (\xx,-0.04) node[below,font=\scriptsize] {\xx};}
% ---- derivative bells (right axis, scale 0.27273) ----
\begin{scope}[yscale=0.27273]
% 268K solid
\draw[thick] plot[smooth] coordinates {(-85,1.0386) (-75,1.5595) (-65,2.3099) (-55,3.3530) (-45,4.7267) (-35,6.3945) (-25,8.1841) (-15,9.7592) (-5,10.6992) (0,10.8251) (5,10.6992) (15,9.7592) (25,8.1841) (35,6.3945) (45,4.7267) (55,3.3530) (65,2.3099) (75,1.5595) (85,1.0386)};
% 298K dashed
\draw[thick,dashed] plot[smooth] coordinates {(-85,1.3235) (-75,1.8898) (-65,2.6585) (-55,3.6627) (-45,4.9035) (-35,6.3180) (-25,7.7495) (-15,8.9500) (-5,9.6436) (0,9.7353) (5,9.6436) (15,8.9500) (25,7.7495) (35,6.3180) (45,4.9035) (55,3.6627) (65,2.6585) (75,1.8898) (85,1.3235)};
% 328K dotted
\draw[thick,dotted] plot[smooth] coordinates {(-85,1.5878) (-75,2.1740) (-65,2.9309) (-55,3.8697) (-45,4.9708) (-35,6.1642) (-25,7.3179) (-15,8.2503) (-5,8.7761) (0,8.8449) (5,8.7761) (15,8.2503) (25,7.3179) (35,6.1642) (45,4.9708) (55,3.8697) (65,2.9309) (75,2.1740) (85,1.5878)};
% peak dots
\fill (0,10.8251) circle (2.2pt); \fill (0,9.7353) circle (2.2pt); \fill (0,8.8449) circle (2.2pt);
% FWHM of 298K bell: half-max 4.8677 /V at (V-U)=+-45.27 mV
\draw[<->,black!55] (-45.27,4.8677) -- (45.27,4.8677);
\node[fill=white,inner sep=1pt,font=\scriptsize,black!55] at (0,4.8677) {FWHM $3.53w{=}90.5$ mV (298 K)};
\end{scope}
% ---- representative logistic xi_eq (298K, left axis, scale 3.0), gray ----
\begin{scope}[yscale=3.0]
\draw[gray] plot[smooth] coordinates {(-85,0.0352) (-70,0.0615) (-55,0.1051) (-40,0.1740) (-25,0.2742) (-10,0.4039) (0,0.5000) (10,0.5961) (25,0.7258) (40,0.8260) (55,0.8949) (70,0.9385) (85,0.9648)};
% center tangent, slope 1/(4w): (-40,0.1106)->(40,0.8894)
\draw[densely dotted] (-40,0.1106) -- (40,0.8894);
\fill[gray] (0,0.5) circle (2.2pt);
\node[gray,font=\scriptsize,anchor=west] at (52,0.78) {$\xi_\eq$ (298 K)};
\node[font=\scriptsize,anchor=north west,align=left] at (10,0.47) {tangent slope\\$=1/(4w)$};
\end{scope}
% ---- legend (quantitative), above the plot to clear the bells ----
\node[anchor=south west,font=\scriptsize,align=left] at (-72,3.35)
  {268 K (solid): $w{=}23.1$ mV, peak $10.83$/V\\
   298 K (dashed): $w{=}25.7$ mV, peak $9.74$/V\\
   328 K (dotted): $w{=}28.3$ mV, peak $8.85$/V};
\end{tikzpicture}
\caption{평형 진행률 $\xi_\eq$ 와 그 미분 --- 폭 $w_j{=}n_jRT/F$ ($n_j{=}1$)의 온도 의존(좌표는 식 그대로의
수치 평가). 오른축 종(bell)이 $\dd\xi_\eq/\dd V=\xi_\eq(1-\xi_\eq)/w$(식~\eqref{eq:belliden})로 세 온도
$T{=}268/298/328$ K 에서 폭 $w{=}23.1/25.7/28.3$ mV, 정점 $1/(4w){=}10.83/9.74/8.85$ /V (검은 점)를 갖는다
--- 온도가 오르면 종이 낮고 넓어진다(면적 보존). 왼축 회색 $S$-곡선이 대응 logistic(식~\eqref{eq:xieq}, 298 K)이고,
중심 접선의 기울기가 정확히 $1/(4w)$ 라 $w$ 의 기하적 의미를 준다. FWHM $=2\ln(3{+}2\sqrt2)\,w=3.53\,w=90.5$ mV
(298 K). 종은 방향 불변(양수)이다.}
\label{fig:logistic}
\end{figure}

\subsection{같은 결과의 분포 관점 --- $\xi_\eq$ 는 평형 점유 분포의 여집합}\label{sec:dist}
앞에서 $\xi_\eq$ 는 두 경로로 나왔다 --- (i) kinetic: 속도식 정지점에서 detailed balance~\eqref{eq:db}(식~\eqref{eq:logisticsolve}),
(ii) thermo: 자유에너지 최소화 $g_j'(\xi)=sF(V_\eq-U)$(식~\eqref{eq:Veq}). 이 둘은 사실 \emph{하나의 분포}의 두 측면이다.
이 소절은 결과식$\cdot$부호$\cdot$모델은 그대로 두고 $\xi_\eq$ 의 통계역학적 정체만 밝혀,
Chapter 2 \S\ref{sec:lco-electronic} 의 LCO 전자 엔트로피와 같은 언어로 두 곳을 묶는다.

\textbf{(a)(b) --- Part 0 회수.} 단일 자리 분포는 Part 0 이 이미 닫았다 --- grand-canonical
분배함수 $\Xi_1=1+e^{-\beta\Delta\mu}$(식~\eqref{eq:partfn}, 내부 자유도 $q(T)$ 흡수형)와 평균 점유
$\langle n\rangle=1/(1+e^{+\beta\Delta\mu})$(식~\eqref{eq:fermifn},
$\Delta\mu\equiv\tilde\varepsilon-\mu_\mathrm{Li}$)를 받는다. \textbf{(c) 같은 식 --- logistic 과의 일치.}
전기화학 구동력이 자리 점유 비용을 정한다 --- 전이 중심 기준(자리 자유에너지의 몰 환산 $N_A\tilde\varepsilon=\mu^0$)과
식~\eqref{eq:eqcond} 의 $\mu_\mathrm{Li}=\mu^0-sF(V-U)$ 를 정의에 넣으면 1몰 기준으로
$\Delta\mu=+sF(V-U_j)$, 곧 지수는 몰당 형태로
$\Delta\mu/(RT)=+s(V-U_j)/w_j$($w_j\equiv RT/F$, 이상 극한)다($\beta$ 의 자리당 $k_BT$ 가 몰당 $RT$ 로,
자리당 전하 $e$ 가 몰당 $F$ 로 한꺼번에 환산 --- 비는 불변). 따라서 식~\eqref{eq:fermifn} 은
$\langle n\rangle=1/(1+e^{+s(V-U_j)/w_j})=\theta_\eq$ --- $V$ 를 올리면 점유가 내려가는(탈리튬화) 자리 점유
그 자체이고, 본론의 진행률은 그 여집합 $\xi_\eq=1-\langle n\rangle=1/(1+e^{-s(V-U_j)/w_j})$ 로
식~\eqref{eq:logisticsolve}$\cdot$\eqref{eq:xieq} 와 일치한다(Part T 의 식~\eqref{eq:logistic}과 부호까지 동일).
곧 \emph{$\xi_\eq$ 는 격자기체 평형 점유 확률 분포의 여집합}이며, 점유 $\theta_\eq$ 와 진행률 $\xi_\eq$ 를
가르는 한 번의 부호 차이는 방향 부호가 아니라 \emph{같은 분포를 읽는 두 좌표}(여집합 교환
$\xi=1-\theta$)의 차이다.

\textbf{(d) 두 경로의 통합.} kinetic(정$\cdot$역 플럭스 균형 점유, 식~\eqref{eq:db})$\cdot$thermo 두 경로
모두 \emph{같은 grand-canonical 평형 점유 분포}(식~\eqref{eq:fermifn})에 도달한다. $\Omega=0$(이상 격자기체)이면 분포가 정확히 Fermi-함수형이고, $\Omega>0$
이면 평균장 자리 비용 $\Delta\mu$ 에 $-\Omega(1-2\xi)$ 가 더해져(식~\eqref{eq:mu}) 자기무모순 분포가 된다 --- 그
비단조성이 \S\ref{sec:hys} 의 히스이고, 그 분포가 접히는 점(자유에너지 $g$ 의 변곡, spinodal)이 식~\eqref{eq:spinodal} 이다. 곧 ``점유 분포''
한 언어가 평형 중심$\cdot$폭$\cdot$히스$\cdot$logistic 을 한 줄에 꿴다.

\begin{keybox}
\textbf{분포 다리.} $\xi_\eq$ 가 평형 점유 확률 분포(식~\eqref{eq:fermifn})라는 것이 Chapter 2 \S\ref{sec:lco-electronic}
의 LCO \emph{전자} 엔트로피(전도 전자 Fermi--Dirac 분포 $f(E)=1/(1+e^{\beta(E-E_F)})$)와 \emph{동형}이다 ---
흑연의 \emph{리튬 자리} 점유와 LCO 의 \emph{전자 준위} 점유가 같은 통계로 닫혀, 이 다리가 LCO 전자 엔트로피
절을 흑연 forward 틀 안으로 편입시킨다.
\end{keybox}

% 자산: [A-101] [A-102] [A-103] [A-104] [A-105] [A-106] [A-107] [A-108] [A-109] [A-110] [A-111] [A-112] [A-113] [A-114] [A-115] [A-116] [A-117] [A-118] [A-119] [A-120] [A-121] [A-122]
